\documentclass{article}

% NeurIPS-inspired formatting (self-contained, no template required)
\usepackage[margin=1in]{geometry}
\usepackage{times}
\usepackage{amsmath,amssymb}
\usepackage{graphicx}
\usepackage{booktabs}
\usepackage{hyperref}
\usepackage{microtype}
\usepackage{caption}
\usepackage{subcaption}
\usepackage[numbers]{natbib}
\usepackage{setspace}

\title{\textbf{Stable Diffusion Parameter Explorer:\\
A Systematic Empirical Study of Inference Hyperparameters}}

\author{Jae Kim\\
COGS 185 --- Final Project\\
University of California, San Diego}

\date{March 2026}

\begin{document}
\maketitle

% ─────────────────────────────────────────────────────────────────
\begin{abstract}
Stable Diffusion has become a widely deployed framework for text-to-image synthesis,
yet practitioners face a poorly characterized inference-time hyperparameter space.
We present a systematic empirical study of four key parameters---classifier-free
guidance (CFG) scale, denoising step count, noise scheduler, and model checkpoint---
across 162 generated images spanning 3 prompts and 3 random seeds per condition.
Image quality is quantified via CLIP cosine similarity (ViT-B/32)
normalized to $[0, 1]$.
Our main findings are: (1)~CFG scale has the largest measurable impact on alignment,
with the steepest gain occurring between CFG~$= 1$ and CFG~$= 3$ ($\Delta{=}+0.0131$),
saturating near CFG~$= 15$ (mean CLIP~$= 0.6637$);
(2)~denoising steps plateau sharply---10~steps achieves 99.7\% of the maximum
CLIP score at one-fifth the compute cost of 50~steps;
(3)~scheduler choice produces negligible quality differences (range $\leq 0.0022$),
with 98.7\% of CLIP variance attributable to prompt/seed rather than scheduler;
(4)~SD v1.4 and SD v1.5 are statistically indistinguishable by this metric
($\Delta{=}0.0002$).
We accompany this study with an open-source Streamlit application for interactive
parameter exploration and a reproducible batch-runner codebase.
\end{abstract}

% ─────────────────────────────────────────────────────────────────
\section{Introduction}

Latent diffusion models~\citep{rombach2022ldm}, and Stable Diffusion in particular,
have transformed text-to-image generation by combining the expressive power of
diffusion-based denoising~\citep{ho2020ddpm} with computationally efficient
latent-space generation.
Despite widespread deployment, there is surprisingly little rigorous empirical
guidance on the \emph{inference-time} parameter settings that practitioners
must choose: the classifier-free guidance (CFG) scale~\citep{ho2022cfg},
the number of denoising steps, the noise scheduler, and the model checkpoint.

Most published work on Stable Diffusion focuses on training objectives, architectures,
or conditioning mechanisms.
Systematic analyses of inference hyperparameters are largely absent from the
academic literature, leaving users to rely on informal community heuristics
(e.g., ``CFG~7.5, 20~steps'' as defaults) without quantitative justification.
This gap motivates the present study.

We make the following contributions:
\begin{itemize}
\item A controlled empirical study of four inference hyperparameters---CFG scale,
  step count, scheduler, and model---using CLIP alignment as a unified,
  reproducible evaluation metric.
\item A Pareto-efficiency analysis of the step-count/quality trade-off,
  demonstrating that the community default of 20~steps is well-placed but
  that 10~steps offers a 5$\times$ efficiency advantage with less than 0.003
  CLIP degradation.
\item An open-source batch experiment runner and interactive Streamlit application
  that allows real-time parameter exploration with live CLIP scoring.
\item Reproducible experimental infrastructure, with cross-experiment DDIM
  consistency verified to floating-point precision (9/9 bitwise matches
  across shared conditions).
\end{itemize}

This work was developed using AI-assisted tooling (Claude Code), allowing rapid
construction of a modular experimental codebase---a methodology consistent with
the course's Option~6 (vibe coding) pathway.

% ─────────────────────────────────────────────────────────────────
\section{Method}

\subsection{Model and Experimental Setup}

We use Stable Diffusion v1.5 (\texttt{runwayml/stable-diffusion-v1-5}) as the
primary model, implemented via the \texttt{diffusers} library~\citep{platen2022diffusers}.
All experiments were run on an NVIDIA GeForce GTX~1080 (8.6~GB VRAM) with
\texttt{torch.float16} precision.

Each experiment fixes all parameters at a baseline and varies exactly one
(the single-variable protocol).
Baseline settings are: CFG~$= 7.5$, 20~denoising steps, DDIM scheduler,
and SD v1.5.
We use three prompts representing distinct semantic domains (nature scene,
urban fantasy, wildlife) and three random seeds (42, 123, 999) per condition,
yielding $n = 9$ observations per condition.

\subsection{Evaluation Metric}

We evaluate image-text alignment using CLIP~\citep{radford2021clip}
cosine similarity via the open-source \texttt{open\_clip} library~\citep{ilharco2021openclip}
with the ViT-B/32 architecture pretrained on OpenAI's dataset.
Raw cosine similarities lie in $[-1, 1]$; we normalize to $[0, 1]$ as:
\begin{equation}
  \text{CLIP Score} = \frac{\text{cos\_sim}(I, T) + 1}{2}
  \label{eq:clip}
\end{equation}
where $I$ is the image embedding and $T$ is the text embedding.
We note that this metric measures \emph{prompt-image alignment}, not
perceptual quality in isolation; images that are photorealistic but
misaligned to the prompt would score lower than simpler aligned outputs.

\subsection{Schedulers}

We compare four samplers: DDIM~\citep{song2021ddim},
PNDM~\citep{liu2022pndm}, DPM++2M~\citep{lu2022dpm}, and
Euler discrete.
All schedulers are initialized from the same model configuration.

\subsection{Interactive Application}

The accompanying Streamlit app (\texttt{app.py}) provides real-time
parameter exploration with live CLIP scoring.
It uses \texttt{@st.cache\_resource} to avoid pipeline reloads on parameter changes,
a batch sweep mode for side-by-side CFG comparisons, and a download button for
saving generated images.
The app extends the paper's experiments by additionally supporting arbitrary prompts,
negative prompts, and on-demand single-image generation.

% ─────────────────────────────────────────────────────────────────
\section{Experiments}

Table~\ref{tab:summary} summarizes all four experiments.

\begin{table}[h]
\centering
\caption{Summary of all experiments. CLIP range is the spread from worst to best
  condition mean; a larger range indicates greater parameter sensitivity.}
\label{tab:summary}
\begin{tabular}{llrrr}
\toprule
Experiment & Variable & Conditions & CLIP Range & Best CLIP \\
\midrule
Exp~1 & CFG Scale      & 7  & 0.0199 & 0.6637 (CFG=15) \\
Exp~2 & Steps          & 5  & 0.0020 & 0.6619 (30 \& 50 steps) \\
Exp~3 & Scheduler      & 4  & 0.0022 & 0.6612 (Euler) \\
Exp~4 & Model (v1.x)   & 2  & 0.0002 & 0.6607 (SD v1.5) \\
\bottomrule
\end{tabular}
\end{table}

\subsection{Experiment 1: Classifier-Free Guidance Scale}

\begin{figure}[h]
  \centering
  \includegraphics[width=0.65\textwidth]{results/fig_cfg_scale.png}
  \caption{Mean CLIP score vs.\ CFG scale (averaged over 3 prompts and 3 seeds).
    The saturation shape is broadly consistent, though sensitivity magnitude varies
    by prompt (see text).}
  \label{fig:cfg}
\end{figure}

We swept CFG~$\in \{1, 3, 5, 7, 10, 15, 20\}$ with DDIM and 20~steps fixed.
Figure~\ref{fig:cfg} and Table~\ref{tab:cfg} show a clear nonlinear relationship
(Pearson $r = 0.48$): alignment rises steeply from CFG~$= 1$ to CFG~$= 3$
($\Delta{=}+0.0131$, representing 66\% of the total achievable gain),
then plateaus, and reverses slightly at CFG~$= 20$ ($-0.0012$ vs.\ CFG~$= 15$).

\begin{table}[h]
\centering
\caption{CFG scale experiment results. Mean $\pm$ std over 9 images per condition.}
\label{tab:cfg}
\begin{tabular}{rcc}
\toprule
CFG Scale & Mean CLIP $\pm$ Std & Time (s) \\
\midrule
1  & $0.6437 \pm 0.0074$ & 5.89 \\
3  & $0.6568 \pm 0.0066$ & 9.05 \\
5  & $0.6601 \pm 0.0068$ & 9.08 \\
7  & $0.6607 \pm 0.0063$ & 9.37 \\
10 & $0.6615 \pm 0.0073$ & 9.22 \\
\textbf{15} & $\mathbf{0.6637 \pm 0.0068}$ & 9.60 \\
20 & $0.6624 \pm 0.0061$ & 9.49 \\
\bottomrule
\end{tabular}
\end{table}

The effect is prompt-dependent: the ``futuristic city'' prompt shows a gain of
$+0.0247$ from CFG~$= 1$ to CFG~$= 15$, while the ``snowy cabin'' prompt gains
only $+0.0125$---a nearly 2$\times$ difference in sensitivity.
This suggests that prompts with more abstract or compositional elements
(\emph{flying cars}, \emph{city skyline}) benefit more from stronger guidance.

Notably, higher CFG also reduces output variability: the coefficient of variation
falls from 1.14\% at CFG~$= 1$ to 0.91\% at CFG~$= 20$, indicating that stronger
guidance not only improves alignment but stabilizes generation.

\textbf{Recommendation:} CFG in the range 7--12 offers near-peak alignment
without over-saturation risk. CFG~$= 15$ is optimal by CLIP but may introduce
visible artifacts in practice.

\subsection{Experiment 2: Denoising Step Count}

\begin{figure}[h]
  \centering
  \includegraphics[width=0.9\textwidth]{results/fig_steps.png}
  \caption{Left: CLIP score vs.\ denoising steps. Right: generation time vs.\ steps.
    Quality plateaus while compute scales nearly linearly.}
  \label{fig:steps}
\end{figure}

We varied step count $\in \{10, 15, 20, 30, 50\}$ with DDIM and CFG~$= 7.5$ fixed.
The CLIP-score range across all conditions is only 0.0020---far smaller than the
CFG effect---while generation time scales from 4.72~s (10~steps) to 23.57~s
(50~steps), a 5$\times$ cost increase.

\begin{table}[h]
\centering
\caption{Steps efficiency analysis. ``\%CLIP$_\text{max}$'' is the fraction of
  the maximum CLIP score achieved. ``CLIP/s'' is the efficiency ratio.}
\label{tab:steps}
\begin{tabular}{rrrrr}
\toprule
Steps & Mean CLIP & Time (s) & CLIP/s & \%CLIP$_\text{max}$ \\
\midrule
\textbf{10} & 0.6599 & 4.72 & \textbf{0.1397} & 99.7\% \\
15 & 0.6611 & 6.81 & 0.0971 & 99.9\% \\
20 & 0.6607 & 8.96 & 0.0738 & 99.8\% \\
30 & 0.6619 & 13.44 & 0.0492 & 100.0\% \\
50 & 0.6619 & 23.57 & 0.0281 & 100.0\% \\
\bottomrule
\end{tabular}
\end{table}

The efficiency ratio (CLIP per second) collapses 5.0$\times$ from
10~steps (0.1397) to 50~steps (0.0281).
Critically, steps 30 and 50 achieve \emph{identical} mean CLIP to four decimal
places (0.6619), yet 50~steps costs 10.1~additional seconds per image---10~s
of pure waste given fixed quality.

\textbf{Recommendation:} 15--20 steps provides near-maximum CLIP quality
at 29--38\% of the 50-step compute budget.
10~steps is appropriate for rapid prototyping with less than 0.003 CLIP penalty.

\subsection{Experiment 3: Noise Scheduler}

\begin{figure}[h]
  \centering
  \includegraphics[width=0.55\textwidth]{results/fig_scheduler.png}
  \caption{Mean CLIP score by scheduler. The total spread of 0.0022 is
    smaller than the within-condition standard deviation for any single scheduler.}
  \label{fig:sched}
\end{figure}

We compared DDIM, PNDM, DPM++2M, and Euler at 20~steps and CFG~$= 7.5$.
The total CLIP range across all four schedulers is 0.0022---smaller than the
within-condition standard deviation of any single scheduler ($\approx 0.007$).

A variance decomposition of the 36 CLIP scores in this experiment reveals that
98.7\% of total variance is \emph{within-condition} (attributable to prompt/seed),
and only 1.3\% is \emph{between-scheduler}.
This is strong evidence that scheduler choice is a second-order factor for
CLIP alignment quality.

\begin{table}[h]
\centering
\caption{Scheduler comparison. Generation times are inflated by pipeline-reload
  overhead (see note) and do not reflect true per-sample inference time.}
\label{tab:sched}
\begin{tabular}{lcc}
\toprule
Scheduler & Mean CLIP $\pm$ Std & CV (\%) \\
\midrule
DDIM     & $0.6607 \pm 0.0077$ & 1.16 \\
PNDM     & $0.6590 \pm 0.0067$ & 1.01 \\
DPM++2M  & $0.6609 \pm 0.0080$ & 1.20 \\
Euler    & $0.6612 \pm 0.0076$ & 1.15 \\
\bottomrule
\end{tabular}
\end{table}

\noindent\textbf{Implementation note:} Generation times in this experiment
reflect full pipeline initialization overhead, as the batch runner reloads
the model for each scheduler change.
True per-sample inference time at 20~steps is approximately 9~s for all
schedulers (consistent with Experiment~2's DDIM baseline).
This is a known artifact of the experimental codebase and does not affect the
quality measurements.

Cross-experiment consistency: all 9 DDIM samples shared between
Experiment~2 (steps~$= 20$) and Experiment~3 match to floating-point precision,
confirming full reproducibility of the experimental framework.

\textbf{Recommendation:} Scheduler choice should be driven by factors other than
CLIP alignment (e.g., determinism preferences or compatibility with specific
LoRA/ControlNet pipelines). DPM++2M is often favored in practice for its
convergence properties at low step counts.

\subsection{Experiment 4: Model Checkpoint (SD v1.4 vs.\ SD v1.5)}

\begin{figure}[h]
  \centering
  \includegraphics[width=0.5\textwidth]{results/fig_model_comparison.png}
  \caption{CLIP score comparison between SD v1.4 and SD v1.5. The difference
    is negligible at the aggregate level but shows a prompt-dependent interaction.}
  \label{fig:model}
\end{figure}

We compared CompVis SD v1.4 and Runway SD v1.5 under identical conditions
(20~steps, CFG~$= 7.5$, DDIM).
The aggregate CLIP difference is 0.0002, which is indistinguishable from noise.
However, a per-prompt breakdown (Table~\ref{tab:model}) reveals a model-prompt
interaction: v1.5 outperforms v1.4 on the ``futuristic city'' prompt
($\Delta{=}+0.0036$) and the ``lion'' prompt ($\Delta{=}+0.0042$),
while v1.4 outperforms on the ``snowy cabin'' prompt ($\Delta{=}-0.0072$).

\begin{table}[h]
\centering
\caption{Model $\times$ prompt CLIP scores (mean over 3 seeds).
  $\Delta$ is v1.5 $-$ v1.4; positive favors v1.5.}
\label{tab:model}
\begin{tabular}{lccc}
\toprule
Prompt & SD v1.4 & SD v1.5 & $\Delta$ \\
\midrule
Cozy cabin in snowy forest      & 0.6598 & 0.6525 & $-0.0072$ \\
Futuristic city with flying cars & 0.6596 & 0.6632 & $+0.0036$ \\
Majestic lion at sunset         & 0.6622 & 0.6664 & $+0.0042$ \\
\midrule
\textbf{Overall}                & 0.6605 & 0.6607 & $+0.0002$ \\
\bottomrule
\end{tabular}
\end{table}

We caution that with $n = 3$ per cell, the interaction pattern is exploratory
and should not be generalized without replication.
The finding is nonetheless consistent with reported improvements in v1.5's
human-evaluation scores for photorealistic subjects (animals, people) relative
to abstract or landscape prompts~\citep{runway2022sd15}.

\textbf{Recommendation:} Within the v1.x checkpoint family, model version is
not a meaningful quality lever by CLIP alignment. Upgrade to v2.x or SDXL if
architectural improvements are desired.

\subsection{Supplementary Study: Prompt Complexity (Qualitative)}

\begin{figure}[h]
  \centering
  \includegraphics[width=0.6\textwidth]{results/fig_prompt_complexity.png}
  \caption{CLIP score as a function of prompt complexity level.
    The trend is broadly increasing but non-monotonic (see text).}
  \label{fig:prompt}
\end{figure}

As a supplementary experiment (results not included in the main CSV),
we compared four levels of prompt specificity for a lion subject:
\emph{Minimal} (``a lion''), \emph{Short} (``a majestic lion in the savanna''),
\emph{Detailed} (adding ``golden hour, photorealistic''), and
\emph{Very Detailed} (adding resolution, fur texture, and lighting descriptors).
Figure~\ref{fig:prompt} shows a broadly increasing trend in CLIP score with prompt
specificity, though not strictly monotonic: the \emph{Short} prompt (0.6511) scores
\emph{lower} than \emph{Minimal} (0.6529) before rising through \emph{Detailed}
(0.6628) and \emph{Very Detailed} (0.6685).
This Minimal$\to$Short dip is consistent with CLIP's behavior when a short prompt
adds context words (e.g., ``savanna'') that are not fully reflected in the generated
image, briefly diluting alignment before richer descriptors restore and strengthen it.

% ─────────────────────────────────────────────────────────────────
\section{Discussion}

Our results support a practical decision hierarchy for SD inference parameters,
ordered by practical importance:

\begin{enumerate}
\item \textbf{CFG scale} ($\Delta = 0.0199$): The most impactful lever.
  Low values ($\leq 3$) severely degrade alignment.
  The range 7--12 is robust; CFG~$= 15$ is optimal but risks saturation.
\item \textbf{Step count} ($\Delta = 0.0020$): Moderate impact on quality,
  large impact on efficiency.
  The community default of 20~steps is well-calibrated;
  10~steps is viable for rapid workflows.
\item \textbf{Scheduler} ($\Delta = 0.0022$): Negligible quality impact.
  98.7\% of variance is prompt/seed-driven.
  Choose for architectural compatibility rather than quality.
\item \textbf{Model checkpoint} (v1.x, $\Delta = 0.0002$): No measurable
  quality difference within the v1.x family.
\end{enumerate}

\noindent\textbf{Limitations.}
CLIP score measures text alignment, not perceptual quality or aesthetic appeal;
the narrow range of scores (0.626--0.675) means differences may be imperceptible
to human viewers.
Our study uses only 3 prompts, limiting generalizability to more complex,
compositional, or abstract text descriptions.
With $n = 9$ per condition, we lack statistical power for formal significance
testing on small effects ($\Delta < 0.005$).
The scheduler timing measurements are confounded by model-reload overhead
and cannot be used to compare inference speed.

% ─────────────────────────────────────────────────────────────────
\section{Conclusion}

We presented a systematic empirical characterization of inference-time
hyperparameters in Stable Diffusion v1.x, generating 162 images across 4
experiments and evaluating alignment via normalized CLIP similarity.
The results show that CFG scale is the dominant quality lever, denoising steps
exhibit strong diminishing returns above 15--20, and both scheduler choice and
v1.x checkpoint version are essentially irrelevant to CLIP alignment.

The practical upshot for SD users: \emph{set CFG~$= 7$--$12$, use 15--20~steps,
and do not spend time choosing schedulers for quality reasons}.

The accompanying Streamlit application makes these findings interactive and
reproducible, and the open-source codebase (\texttt{experiments.py},
\texttt{evaluate.py}) enables extension to new prompts, models, or metrics.

\section*{Bonus Points}
This project makes several contributions that merit consideration for bonus points:

\begin{itemize}
\item \textbf{Efficiency Pareto analysis (novel contribution):}
  The CLIP/s efficiency metric and the 5$\times$ degradation finding from
  10 to 50 steps is, to our knowledge, not previously published in this form.
  It provides a directly actionable decision rule for practitioners.
\item \textbf{AI-assisted development (Option 6):}
  The full experimental infrastructure (batch runner, evaluator, Streamlit app)
  was designed and implemented using Claude Code, demonstrating effective
  vibe-coding methodology for scientific computing.
  The resulting modular architecture (decoupled \texttt{evaluate.py}/\texttt{experiments.py})
  reflects best practices achievable through AI-assisted iteration.
\item \textbf{Reproducibility:} The cross-experiment DDIM consistency check
  (9/9 bitwise matches) provides a strong internal validity guarantee not commonly
  seen in student projects.
\item \textbf{Interactive system:} The Streamlit app extends the paper with
  a live parameter exploration tool, supporting negative prompts and batch sweeps
  beyond the scope of the written experiments.
\end{itemize}

% ─────────────────────────────────────────────────────────────────
\bibliographystyle{plainnat}
\bibliography{references}

\end{document}
